\documentclass[10pt]{beamer}
\usetheme{Warsaw}
\usepackage[T1]{fontenc}
\usepackage[utf8]{inputenc}
\usepackage{polski}
\usepackage{amsmath}
\usepackage{amsfonts}
\usepackage{amssymb}
\usepackage{amsbsy}
\usepackage{graphicx}
\usepackage{listings}
\usepackage{xcolor}

\definecolor{codegreen}{rgb}{0,0.6,0}
\definecolor{codegray}{rgb}{0.5,0.5,0.5}
\definecolor{codepurple}{rgb}{0.58,0,0.82}
\definecolor{backcolour}{rgb}{0.95,0.95,0.92}

\lstset{
    backgroundcolor=\color{backcolour},   
    commentstyle=\color{codegreen},
    keywordstyle=\color{magenta},
    numberstyle=\tiny\color{codegray},
    stringstyle=\color{codepurple},
    basicstyle=\ttfamily\footnotesize,
    breakatwhitespace=false,         
    breaklines=true,                 
    captionpos=b,                    
    keepspaces=true,                 
    numbers=left,                    
    numbersep=5pt,                  
    showspaces=false,                
    showstringspaces=false,
    showtabs=false,                  
    tabsize=2,
    literate={ą}{{\k{a}}}1 {ć}{{\'c}}1 {ę}{{\k{e}}}1 {ł}{{\l{}}}1 {ń}{{\'n}}1 {ó}{{\'o}}1 {ś}{{\'s}}1 {ź}{{\'z}}1 {ż}{{\.z}}1 {Ą}{{\k{A}}}1 {Ć}{{\'C}}1 {Ę}{{\k{E}}}1 {Ł}{{\L{}}}1 {Ń}{{\'N}}1 {Ó}{{\'O}}1 {Ś}{{\'S}}1 {Ź}{{\'Z}}1 {Ż}{{\.Z}}1
}

\author{Tomasz Solarz}
\title[Azure CLI \hspace{2em}\insertframenumber/\inserttotalframenumber]{
	Automatyzacja i zarządzanie zasobami w chmurze za pomocą Azure CLI
}

\logo{\includegraphics[height=1cm]{media/logoMS.png}} 
\institute{
	\includegraphics[height=1cm]{media/logo.jpg}\\
	Wydział Matematyki Stosowanej\\
	Katedra Metod Matematycznych w Technice i Informatyce (RMS3)
}
\date{\today}  

\begin{document}
\begin{frame}
	\titlepage
\end{frame}

\section*{Agenda}
\begin{frame}
    \frametitle{Agenda}
	\tableofcontents
\end{frame}

\section{Wprowadzenie}
\begin{frame}
    \frametitle{Czym jest Azure CLI?}
    \begin{block}{Definicja}
        Azure Command-Line Interface (CLI) to wieloplatformowe narzędzie wiersza poleceń służące do łączenia się z platformą Azure i wykonywania poleceń administracyjnych w jej zasobach.
    \end{block}
    \begin{itemize}
        \item Pozwala na automatyzację zadań poprzez skrypty.
        \item Dostępne na Windows (PowerShell/CMD), macOS, Linux.
        \item Działa również w Azure Cloud Shell bezpośrednio w przeglądarce.
    \end{itemize}
\end{frame}

\section{Instalacja i Logowanie}
\begin{frame}[fragile]
    \frametitle{Pierwsze kroki}
    \begin{block}{Instalacja (Windows)}
        Najprościej za pomocą instalatora MSI lub polecenia:
        \begin{lstlisting}[language=bash]
winget install -e --id Microsoft.AzureCLI
        \end{lstlisting}
    \end{block}
    \begin{block}{Logowanie}
        Aby połączyć się z kontem Azure, należy użyć:
        \begin{lstlisting}[language=bash]
az login
        \end{lstlisting}
        Otworzy to okno przeglądarki w celu uwierzytelnienia.
    \end{block}
\end{frame}

\section{Podstawowe polecenia}
\begin{frame}[fragile]
    \frametitle{Zarządzanie Grupami Zasobów}
    Grupa zasobów (Resource Group) to logiczny kontener na zasoby.
    \begin{lstlisting}[language=bash]
# Tworzenie grupy zasobów
az group create --name MojaGrupa --location westeurope

# Lista grup zasobów
az group list --output table

# Usuwanie grupy (i wszystkich jej zasobów)
az group delete --name MojaGrupa --yes --no-wait
    \end{lstlisting}
\end{frame}

\section{Przykłady użycia}
\subsection{Maszyny Wirtualne}
\begin{frame}[fragile]
    \frametitle{Tworzenie Maszyny Wirtualnej (VM)}
    Azure CLI pozwala na szybkie powołanie infrastruktury:
    \begin{lstlisting}[language=bash]
az vm create \
  --resource-group MojaGrupa \
  --name MojaMaszyna \
  --image Ubuntu2204 \
  --admin-username azureuser \
  --generate-ssh-keys
    \end{lstlisting}
    \begin{itemize}
        \item Powyższe polecenie tworzy VM wraz z siecią, adresem IP i kluczami SSH.
    \end{itemize}
\end{frame}

\subsection{Web Apps}
\begin{frame}[fragile]
    \frametitle{Wdrażanie Web App}
    Azure CLI ułatwia wdrażanie aplikacji webowych:
    \begin{lstlisting}[language=bash]
# Tworzenie planu usługi App Service
az appservice plan create --name MojPlan --resource-group MojaGrupa --sku B1 --is-linux

# Tworzenie aplikacji webowej (Python)
az webapp create --resource-group MojaGrupa --plan MojPlan --name MojaAplikacja316056 --runtime "PYTHON:3.9"
    \end{lstlisting}
\end{frame}

\section{Zaawansowane funkcje}
\begin{frame}[fragile]
    \frametitle{Zapytania i Formaty Wyjściowe}
    Azure CLI wspiera filtrowanie wyników za pomocą zapytań JMESPath:
    \begin{lstlisting}[language=bash]
# Pobranie tylko adresu IP maszyny wirtualnej
az vm list-ip-addresses --name MojaMaszyna --query "[].virtualMachine.network.publicIpAddresses[0].ipAddress" -o tsv
    \end{lstlisting}
    Dostępne formaty wyjściowe (`--output` / `-o`):
    \begin{itemize}
        \item \textbf{json} (domyślny)
        \item \textbf{table} (czytelna tabela)
        \item \textbf{tsv} (wartości oddzielone tabulatorami)
        \item \textbf{yaml}
    \end{itemize}
\end{frame}

\section{Repozytorium i kod}
\begin{frame}
    \frametitle{Repozytorium projektu}
    \begin{block}{Kod źródłowy i skrypty}
        Przykłady skryptów automatyzujących (`.sh` / `.ps1`) oraz kody źródłowe infrastruktury znajdują się w repozytorium:
    \end{block}
    \begin{center}
        \url{https://github.com/tsolarz/ZSI-AzureCLI-Project}
    \end{center}
    \begin{itemize}
        \item Skrypty tworzące pełne środowisko testowe.
        \item Szablony ARM i pliki Bicep.
        \item Przykłady integracji z CI/CD (GitHub Actions).
    \end{itemize}
\end{frame}

\section{Podsumowanie}
\begin{frame}
    \frametitle{Podsumowanie i wnioski}
    \begin{block}{Zalety korzystania z Azure CLI}
    \begin{itemize}
        \item \textbf{Szybkość:} Znacznie szybsze niż przeklikiwanie się przez Portal Azure.
        \item \textbf{Powtarzalność:} Możliwość zapisania kroków w formie skryptu.
        \item \textbf{Integracja:} Łatwe użycie w potokach DevOps.
        \item \textbf{Wszechstronność:} Dostęp do niemal każdej usługi Azure.
    \end{itemize}
    \end{block}  
\end{frame}

\begin{frame}
	\frametitle{Dziękuję za uwagę!}
	\begin{center}
	\textbf{Tomasz Solarz}\\[1cm]
	\textbf{Dziękuję za uwagę!}\\
    Pytania?
	\end{center}
\end{frame}
\end{document}
